% ============================================================================
% Tests report: conclusions.
% ============================================================================
\section{Conclusiones}\label{sec:conclusiones}

\begin{itemize}
  \item Los tres ensayos (22 robots, modo \texttt{QR}, $\phi \approx 0{,}58$) se siguieron completos:
        \SI{100}{\percent} de posiciones observadas, menos del \SI{0.13}{\percent} con umbral relajado y
        tres anclas manuales en total.
  \item La escala de tiempo $\fr = 3$~cuadros/s queda verificada por la configuración de la cámara, la
        duración del registro y la correlación con el sensor (máxima entre \num{2.9975} y \num{3.0025}).
        El inicio de cada video recortado coincide con el fin del pulso LED, que sirve de sincronización.
  \item El comportamiento es reproducible: rapidez mediana \num{0.58}--\SI{0.70}{\milli\metre\per\second},
        percentil 95 \num{3.5}--\SI{3.6}{\milli\metre\per\second}, $|\omega|$ mediana
        \num{3.0}--\SI{3.8}{\degree\per\second} y desplazamiento cuadrático medio superdifusivo
        ($\alpha \approx 1{,}3$) en los tres.
  \item Corregido el espejado del video, 14 de 22 robots giran en neto a la derecha (horario), como
        indica el modo \texttt{QR}, y el grupo circula débilmente en ese mismo sentido. Los otros 8
        contrarrotan despacio y pasan más tiempo contra la pared, como un disco que rueda por el interior
        del recinto.
  \item El grupo se ordena en tres capas concéntricas separadas por aproximadamente un diámetro.
  \item Los atascos colectivos son la principal diferencia entre ensayos (\SI{17}{\percent},
        \SI{3}{\percent} y \SI{6}{\percent} del tiempo). En dos de los tres ensayos coinciden con la
        máxima fuerza registrada por el sensor de presión: en \Cmp{Lab}{A}, el atasco de
        \SI{382}{\second} coincide con el episodio de saturación del sensor.
\end{itemize}

\paragraph{Próximos pasos.}
\begin{itemize}
  \item Marcar en el video la posición del sensor y relacionar la fuerza con los robots en contacto
        con él.
  \item Repetir con los otros modos (\texttt{R}, \texttt{QZ}) y con distinta cantidad de robots, para ver
        cómo cambian los atascos con $\phi$.
  \item Identificar a cada robot entre videos (por ejemplo, con marcas de color) para saber si los que
        contrarrotan son siempre los mismos o los que quedan contra la pared.
\end{itemize}
